\documentclass[compress]{beamer}
\usepackage[latin1]{inputenc}
\usepackage{cancel}
\usepackage{alltt}
\newdimen\topcrop
\topcrop=10cm  %alternatively 8cm if the pdf inclusions are in letter format
\newdimen\topcropBezier
\topcropBezier=19cm %alternatively 16cm if the inclusions are in letter format

\setbeamertemplate{footline}[frame number]
\title{A single number type for Math education\\in Type Theory}
\author{Yves Bertot\\Thomas Portet}
\date{April 2026}
\mode<presentation>
\begin{document}

\maketitle
\begin{frame}
\frametitle{The context}
\begin{itemize}
\item Point of view: young students in mathematics have a single type of numbers in mind
\item Type theory proposes a variety of type numbers
\begin{itemize}
\item Inductive types come with a builtin capability for computation
\item Inductive types come with a direct notion of reasoning by induction
\item Real numbers live a different part of the world
\end{itemize}
\item Experimenting with using only one type: real numbers
\begin{itemize}
\item Use subsets for natural numbers
\item Recover recursive computation
\item Make {\tt ring} and {\tt field} more comfortable
\end{itemize}
\end{itemize}
\end{frame}
\begin{frame}
\frametitle{Plan}
\begin{itemize}
\item The subset of R that contains natural numbers
\item Recursive computation with real numbers
\item A poor man's calc tactic
\item Deep ring and field tactics
\item Showcase examples
\begin{itemize}
\item Trigonometry up to Vi�te's formula
\item Fibonacci upto a closed formula
\item Binomials as fractions of recursively defined factorials
\end{itemize}
\end{itemize}
\end{frame}
\begin{frame}
\frametitle{Subsets instead of types}
\begin{itemize}
\item Avoid using the type {\tt nat}, but use the concept
\item Natural numbers are definable as a subset, inductively
\item Integers can also be defined inductively, but it is minor
\end{itemize}
\end{frame}
\begin{frame}[fragile]
\frametitle{Inductive predicate in Rocq}
\begin{alltt}
Require Import Reals.
Open Scope R_scope.

Inductive Rnat : R -> Prop :=
  Rnat0 : Rnat 0
| Rnat_succ : forall n, Rnat n -> Rnat (n + 1).
\end{alltt}
Generated induction principle:
\begin{alltt}
nat_ind
     : forall P : R -> Prop,
       P 0 ->
       (forall n : R, Rnat n -> P n -> P (n + 1)) ->
       forall r : R, Rnat r -> P r
\end{alltt}
\end{frame}
\begin{frame}
\frametitle{Ad hoc proofs of membership}
\begin{itemize}
\item When \(n, m \in {\mathbb N}\), \((n + m), (n \times m) \in {\mathbb N}\)
\item When \(n, m \in {\mathbb N}\) and \(m \leq n\), \((n - m) \in {\mathbb N}\) can
also be proved
\item This requires an explicit proof
\item Probably good in a training context for students
\item Similar for division
\end{itemize}
\end{frame}
\begin{frame}[fragile]
\frametitle{Computation}
\begin{itemize}
\item Notations are so that 12 (as a real number) is actually {\tt IZR 13}
\item {\tt IZR~:~Z -> R} is a recursive function
\item The interaction with {\tt Compute} is not pleasant
\end{itemize}
\begin{alltt}
Compute 12.
\textcolor{blue}{ = (R1 + R1) * ((R1 + R1) * (R1 + (R1 + R1)))
     : R}
\end{alltt}
\begin{itemize}
\item We designed a new command called {\tt R\_compute}.
\end{itemize}
\begin{alltt}
R_compute (12 + 13).
\textcolor{blue}{ = 25}
\end{alltt}
\begin{itemize}
\item Does not cover \(\lambda\)-calculus or let-binders
\item Works by using a database of registered functions
\item Can produce a lemma stating the result as an equality
\end{itemize}
\end{frame}
\begin{frame}[fragile]
\frametitle{Recursive definitions, following the {\tt nat} recursive scheme}
\begin{itemize}
\item We provide a command {\tt Recursive} that mimicks recursive definitions
with natural numbers
\item Only ``well-defined'' for elements of {\tt Rnat}
\item The command produces two objects
\begin{itemize}
\item The function of type {\tt R -> R}
\item The proof of the logical statement for that function
\item The function is registered in the database for {\tt R\_compute}
\end{itemize}
\end{itemize}
\begin{small}
\begin{verbatim}
Recursive (def fib such that
             fib 0 = 0 /\
             fib 1 = 1 /\
             forall n : R, Rnat (n - 2) ->
                fib n = fib (n - 2) + fib (n - 1)).
\end{verbatim}
\end{small}
\end{frame}
\begin{frame}
\frametitle{Recursive definitions of functions with several arguments}
Thomas : we need to describe your work here
\end{frame}



\begin{frame}
\frametitle{Minimal set of tactics}
\begin{itemize}
\item {\tt replace}
\begin{itemize}
\item {\tt ring} and {\tt field} for justifications
\item No need to massage formulas step by step through rewriting
\end{itemize}
\item {\tt intros}, {\tt exists}, {\tt split}, {\tt destruct} to handle
  logical connectives (as usual)
\item {\tt rewrite} to handle the behavior of all defined functions
 (and recursors)
\item {\tt unfold} for functions defined by students
\begin{itemize}
\item But we should block unfolding of recursive functions
\end{itemize}
\item {\tt apply} and {\tt lra} to handle all side conditions related to bounds
\item {\tt typeclasses eauto} to prove membership in {\tt Rnat}
\begin{itemize}
\item Explicit handling for subtraction and division
\end{itemize}
\item Possibility to add ad-hoc computing facilities for user-defined
\begin{itemize}
\item Relying on mirror functions computing on inductive {\tt nat} or {\tt Z}
\end{itemize}
\end{itemize}
\end{frame}
\begin{frame}
\frametitle{Demonstration time}
\begin{itemize}
\item A study of factorials and binomial numbers
\begin{itemize}
\item Efficient computation of factorial numbers
\item Proofs relating the two points of view on binomial numbers,
   ratios or recursive definition
\item A proof of the expansion of \((x + y) ^ n\)
\end{itemize}
\item A study the fibonacci sequence
\begin{itemize}
\item \({\cal F}(i) = \frac{\phi ^ i - \psi ^ i}{\phi - \psi}\) (\(\phi\) golden ratio)
\end{itemize}
\end{itemize}
\end{frame}
\begin{frame}
\frametitle{The ``17'' exercise}
\begin{itemize}
\item Prove that there exists an \(n\) larger than 4 such that\\
\[\left(\begin{array}{c}n\\5\end{array}\right) =
  17 \left(\begin{array}{c}n\\4\end{array}\right)\]\\
(suggested by S. Boldo, F. Cl�ment, D. Hamelin, M. Mayero, P. Rousselin)
\item Easy when using the ratio of factorials and eliminating common
sub-expressions on both side of the equality\\
\[\frac{\cancel{n!}}{\cancel{(n - 5)!}\cancel{5!}_5} = 17 \frac{\cancel{n!}}{\cancel{(n-4)!}_{(n-4)}\cancel{4!}}\]
\item They use the type of natural numbers and equation\\
\[\left(\begin{array}{c}n\\p + 1\end{array}\right) \times (p + 1) =
  \left(\begin{array}{c}n\\p\end{array}\right) \times (n - p)\]\\
\end{itemize}
\end{frame}
\end{document}


%%% Local Variables: 
%%% mode: latex
%%% TeX-master: t
%%% End: 
